\DocumentMetadata{
  pdfversion=2.0,
  pdfstandard=ua-2,
  lang=en-US,
  tagging=on,
  tagging-setup={math/setup=mathml-SE,tabsorder=structure}
}
\documentclass[11pt,letterpaper]{article}
\usepackage[margin=0.68in,includefoot]{geometry}
\usepackage{newcomputermodern}
\usepackage{amsmath,amscd,mathtools}
\usepackage{tikz,adjustbox}
\providecommand{\square}{\mdlgwhtsquare}
\usepackage[hidelinks]{hyperref}
\hypersetup{
  pdftitle={Analysis First-Year Exam, May 2022, Part 1},
  pdfauthor={University of Florida Department of Mathematics},
  pdfsubject={Graduate mathematics examination},
  pdfdisplaydoctitle=true,
  bookmarksopen=true
}
\setcounter{secnumdepth}{0}
\setlength{\parindent}{0pt}
\setlength{\parskip}{0.28em}
\setlength{\emergencystretch}{3em}
\setlength{\leftmargini}{2.15em}
\setlength{\leftmarginii}{2.65em}
\setlength{\leftmarginiii}{3em}
\renewcommand{\labelenumi}{\textbf{\theenumi.}}
\pagestyle{empty}
\raggedbottom
\allowdisplaybreaks
\newenvironment{examproblems}[1]{%
  \begin{enumerate}%
  \setcounter{enumi}{#1}%
  \setlength{\topsep}{0.35em}%
  \setlength{\partopsep}{0pt}%
  \setlength{\itemsep}{0.48em}%
  \setlength{\parsep}{0.18em}%
}{\end{enumerate}}
\newenvironment{examsubparts}{%
  \begin{enumerate}%
  \setlength{\topsep}{0.18em}%
  \setlength{\partopsep}{0pt}%
  \setlength{\itemsep}{0.16em}%
  \setlength{\parsep}{0.1em}%
}{\end{enumerate}}
\newenvironment{examinstructions}{%
  \par\begingroup\itshape
}{%
  \par\endgroup\vspace{0.18em}
}
\makeatletter
\renewcommand\section{\@startsection{section}{1}{0pt}%
  {-1.25ex plus -.25ex minus -.1ex}%
  {0.55ex plus .1ex}%
  {\normalfont\large\bfseries}}
\renewcommand{\@maketitle}{%
  \newpage\null\vskip -1em%
  {\footnotesize\bfseries
    \href{https://www.ufl.edu/}{University of Florida}, \href{https://math.ufl.edu/}{Department of Mathematics}\par}%
  \vskip -0.5em%
  \tagstructbegin{tag=Title}%
    {\LARGE\bfseries\href{https://gma.math.ufl.edu/exams/analysis-first-year/}{Analysis First-Year Exam}, May 2022, Part 1\par}%
  \tagstructend%
  \vskip 0.25em%
  \tagmcbegin{artifact}%
    \hrule height 1pt%
  \tagmcend%
  \vskip 1em%
}
\makeatother
\title{Analysis First-Year Exam, May 2022, Part 1}
\author{}
\date{}
\begin{document}
\maketitle
\thispagestyle{empty}
\begin{examinstructions}
Answer FOUR questions in detail. State carefully any results used without proof.
\end{examinstructions}
\begin{examproblems}{0}
\item
Let \((a_n)_{n=0}^{\infty}\) and \((b_n)_{n=0}^{\infty}\) be bounded real sequences. Prove that 
\[
\liminf(a_n+b_n) \geq \liminf a_n+\liminf b_n
\]
 and show by example that the inequality can be strict.
\item
Let the metric space~\(M\) be compact; let \(f:M\to M\) be distance-preserving and let \(a\in M\). Prove that~\(a\) is in the closure of the sequence \(f(a),f(f(a)),\ldots\) that~\(a\) generates by repeated application of~\(f\). Conclude that~\(f\) is surjective.
\item
Determine precisely all real numbers~\(x\) for which the series 
\[
\sum_{n=0}^{\infty}\frac{x^n}{1+x^n}
\]
 is convergent.
\item
Prove that no real-valued function on the unit square \([0,1]\times[0,1]\) can be simultaneously continuous and injective.
\item
Let \(f:(0,\infty)\to\mathbb{R}\) be differentiable; let~\(A\) and~\(L\) be real numbers.
\begin{examsubparts}
\item[\textnormal{(i)}]
Define what is meant by the statement \(\lim_{x\to\infty}f(x)=A\).
\item[\textnormal{(ii)}]
Prove that if \(\lim_{x\to\infty}f(x)=A\) and \(\lim_{x\to\infty}f'(x)=L\), then \(L=0\).
\end{examsubparts}
\end{examproblems}
\end{document}
